\chapter{Ранговые страты частичных изометрий и внешний квадрат моста}

Самосогласованное очищение устранило бесконечную мостовую долину, но
оставило три нулевые страты рангов $1$, $2$ и $3$. Настоящая проверка
проверяет все уже доступные селекторы до введения нового поля: плоскую
меру, энтропию, Real-структуру, топологию проекторного состояния и
архивный double-path механизм Тома IV.

\section{Геометрия трёх нулевых страт}

Обозначим через $\mathcal M_k$ множество вещественных частичных изометрий
$B\in M_3(\mathbb R)$ ранга $k$. Выбор начальной $k$-плоскости, конечной
$k$-плоскости и ортогонального отождествления между ними даёт
\begin{equation}
 \dim\mathcal M_k
 =2k(3-k)+\frac{k(k-1)}2.
\end{equation}
Следовательно,
\begin{equation}
 \dim\mathcal M_1=4,
 \qquad
 \dim\mathcal M_2=5,
 \qquad
 \dim\mathcal M_3=3.
 \label{eq:v6-partial-isometry-dimensions}
\end{equation}
Это стандартная геометрия Штифеля и Грассмана для фиксированного ранга;
см. работу Edelman, Arias и Smith, DOI
\path{10.1137/S0895479895290954}, и
\path{arXiv:1209.0430}.

В канонической точке $B_k=\operatorname{diag}(I_k,0)$ нормальное
пространство состоит из симметричной части активного блока и полного
нулевого блока. Точный спектр гессиана самосогласованного действия равен
\begin{align}
 k=1:&\quad 0^{(4)},\ (4/7)^{(4)},\ (16/7)^{(1)},\nonumber\\
 k=2:&\quad 0^{(5)},\ (2/7)^{(1)},\ (8/7)^{(3)},\nonumber\\
 k=3:&\quad 0^{(3)},\ (16/21)^{(6)}.
 \label{eq:v6-rank-stratum-hessians}
\end{align}
Нулевые кратности точно совпадают с
\eqref{eq:v6-partial-isometry-dimensions}; каждая страта является чистым
минимумом Морса--Ботта.

\section{Плоская мера выбирает ранг два}

Для большого обратного параметра $\beta$ вклад окрестности
$\mathcal M_k$ имеет вид
\begin{equation}
 Z_k(\beta)\sim C_k\,
 \beta^{-(9-\dim\mathcal M_k)/2},
 \qquad C_k>0.
 \label{eq:v6-morse-bott-laplace-rank-weight}
\end{equation}
Это обычная формула Лапласа для многообразия минимумов Морса--Ботта; см.
\path{arXiv:1607.05152} и современную стратифицированную формулировку
\path{arXiv:2608.02626}.

Показатели для рангов $1,2,3$ равны соответственно
\begin{equation}
 -\frac52,\qquad -2,\qquad -3.
\end{equation}
Поэтому при $\beta\to\infty$ плоская мера предпочитает ранг два, а не
ранг один. Это не зависит от конечных объёмных множителей $C_k$.

\begin{theorem}[Флуктуационный селектор ранга два]
Самосогласованное действие вместе с плоской мерой не рождает чистое
состояние. На низкотемпературной асимптотике доминирует пятимерная страта
$\mathcal M_2$, поскольку она имеет наибольшее число мягких направлений.
\end{theorem}

\section{Топология и Real-структура не различают ранги один и два}

Правое редуцированное состояние на нулевой страте имеет вид
\begin{equation}
 R_R(B)=\frac1kP_k.
\end{equation}
Для $k=1$ его орбита есть пространство прямых
$\operatorname{Gr}(1,3)=\mathbb{RP}^2$. Для $k=2$ орбита плоскостей
\begin{equation}
 \operatorname{Gr}(2,3)\simeq\operatorname{Gr}(1,3)
 \simeq\mathbb{RP}^2
\end{equation}
посредством ортогонального дополнения. Обе фазы допускают один и тот же
тип проекторного ежа. Общая гомотопическая классификация дефектов в
упорядоченных средах изложена у Мермина,
\path{doi:10.1103/RevModPhys.51.591}; нематические неоднозначности
подробно разобраны в
\path{doi:10.1103/RevModPhys.84.497}.

Транспонирование $B\mapsto B^T$, реализующее обмен двух Real-ветвей,
сохраняет ранг. Следовательно, ни проекторная топология, ни Real-обмен не
являются селекторами $1$ против $2$.

Строгий класс $\pm15$ Тома V также не решает этот вопрос автоматически.
Он получен после выбора голономной линии и сферически-эквивариантного ежа.
В проекте нет теоремы, отождествляющей $\rank B=k$ с зарядом
$\pm15k$. Использовать примитивность класса для предварительного выбора
$k=1$ означало бы предположить искомый сектор.

\section{Возвращение double-path селектора Тома IV}

В проверке
\path{version4_pati_salam_rank_selector_archaeology_gate.tex}
уже встречался точный ранговый инвариант: разность прямого и перекрёстного
квадратичных путей. Для текущего вещественного моста положим
\begin{equation}
 W_{ia,jb}(B)=B_{ia}B_{jb}-B_{ib}B_{ja}.
\end{equation}
Прямое суммирование даёт
\begin{equation}
 \boxed{
 \|W(B)\|_F^2
 =4\|\bigwedge{}^2B\|_F^2
 =4e_2(B^TB),}
 \label{eq:v6-exterior-square-identity}
\end{equation}
где
\begin{equation}
 e_2(X)=\frac12\bigl((\Tr X)^2-\Tr X^2\bigr).
\end{equation}
Этот функционал неотрицателен и обращается в нуль тогда и только тогда,
когда $\rank B\leq1$. Коэффициент четыре является комбинаторным
результатом антисимметризации, а не непрерывной настройкой.

Однако ненормированный член \eqref{eq:v6-exterior-square-identity}
зависит от общей амплитуды и изменяет уже замороженную скалярную ветвь.
Самосогласованное состояние подсказывает единственную безразмерную форму:
\begin{equation}
 \mathcal W(R)
 =\frac{\|W(B)\|_F^2}{\Tr(B^TB)^2}
 =4e_2(R)
 =2\bigl(1-\Tr R^2\bigr).
 \label{eq:v6-normalized-wedge-state}
\end{equation}
Добавление одной копии $\mathcal W$ к энтропии эквивалентно, с точностью
до константы, уже исследованному функционалу
\begin{equation}
 \Tr(R\log R)-2\left(\Tr R^2-\frac13\right).
 \label{eq:v6-wedge-purity-coefficient}
\end{equation}
Его коэффициент чистоты равен $2>\log4$, то есть он проходит точный порог
перехода предыдущего qutrit-проверки.

Минимум имеет одноосный спектр
\begin{equation}
 R_*=\operatorname{diag}(a,b,b),
 \qquad b=\frac{1-a}{2},
 \qquad
 \log\frac ab=4(a-b),
\end{equation}
где
\begin{equation}
 a=0.95393\ldots,
 \qquad b=0.02303\ldots.
 \label{eq:v6-wedge-ordered-spectrum}
\end{equation}
Все собственные значения положительны, но выделенная собственная линия
имеет орбиту $\mathbb{RP}^2$. Значит, для физики дефекта точный ранг один
не нужен: достаточно полного ранга с расщеплением $1+2$, как уже было
установлено в Томе VI.

\begin{proposition}[Рабочий механизм на уровне потенциала]
Нормированная норма внешнего квадрата с комбинаторным коэффициентом
double-path конструкции переводит изотропное состояние в полноранговую
одноосную фазу $\mathbb{RP}^2$ без непрерывного параметра. Она тем самым
воспроизводит недостающий отрицательный инвариант чистоты на уровне
потенциала.
\end{proposition}

Точная граница остаётся существенной. При любой конечной гладкой связи
энтропийный член $\varepsilon\log\varepsilon$ отталкивает минимум от
границы пространства состояний. Поэтому внешний квадрат выбирает
\emph{почти чистую}, но полноранговую фазу, а не строгий проектор.

\section{Почему механизм ещё не выведен}

Текущий минимальный родитель содержит одну нечётную матрицу $B$ и её
обычный квадрат в кривизне. Он не содержит автоматически
антисимметризованного двухчастичного канала $B\wedge B$. Перенос формулы
из Pati--Salam-графа недопустим без построения прямого и перекрёстного
путей внутри именно обменного соответствия.

\begin{nogocriterion}[Запрет ручного внешнего квадрата]
Функционал \eqref{eq:v6-normalized-wedge-state} нельзя объявлять частью
родительского действия только потому, что он даёт правильную фазу.
Необходимо вывести две квадратичные композиции, их относительный минус,
нормировку на $\Tr(B^TB)^2$ и отсутствие нового независимого веса из
одной Real/градуированной конструкции.
\end{nogocriterion}

\section{Вердикт}

\begin{protocol}
\begin{enumerate}
 \item размерности нулевых страт рангов $1,2,3$:
 \textbf{$4,5,3$};
 \item нормальные гессианы: \textbf{положительны и вычислены точно};
 \item выбор плоской мерой: \textbf{ранг два};
 \item выбор энтропией: \textbf{ранг три};
 \item выбор транспонированием Real-пары: \textbf{нет};
 \item различение рангов $1$ и $2$ топологией состояния:
 \textbf{нет, обе орбиты $\mathbb{RP}^2$};
 \item строгая необходимость точного ранга один для дефекта:
 \textbf{нет, достаточно расщепления $1+2$};
 \item double-path внешний квадрат: \textbf{точный селектор найден};
 \item эквивалентный коэффициент чистоты: \textbf{$2>\log4$};
 \item одноосный полноранговый минимум: \textbf{получен на уровне потенциала};
 \item происхождение внешнего квадрата из текущего родителя:
 \textbf{не доказано};
 \item следующая проверка:
 \textbf{родитель прямого и перекрёстного внешнеквадратных путей};
 \item рождение материи: \textbf{сильно сужено, но не закрыто}.
\end{enumerate}
\end{protocol}

Следующий тест должен построить
\path{version6_exchange_bridge_exterior_square_parent_gate.tex} и
ответить, возникает ли нормированный $B\wedge B$ из одной допустимой
Real-суперсвязности. Если для него требуется новый независимый канал или
ручной коэффициент, механизм остаётся лишь эффективной моделью.

Машинный сертификат сохранён в
\path{s2t_v6_partial_isometry_rank_stratum_selection_gate_results.json}.